\section{What you are preparing to understand}
\label{sec:orientation}

A nearly periodic surface wave can look regular while a small, long-scale
disturbance grows on top of it. Dysthe asks how well an approximate envelope
equation predicts that growth. The cubic nonlinear Schr\"odinger equation,
abbreviated NLS, is the leading model. He retains the next order in wave
steepness and finds that the current induced by a modulated wave changes
the prediction substantially \cite{dysthe}.

The paper is short because it assumes considerable background. Its hardest
steps are the choices of scales, harmonic components, and approximation order
that occur before the linear algebra begins. Your PDE background should make
Laplace's equation, Fourier modes, and the eigenvalue calculation familiar
once their physical roles are identified.

\subsection{A useful reading route}
Read Sections~\ref{sec:waves}--\ref{sec:asymptotics} first. They explain the
wave envelope, the moving free boundary, and why a slow derivative changes
the order of a term. Then study the NLS calculation in
Section~\ref{sec:nls} and the mean-flow calculation in
Section~\ref{sec:mean}. Finally, use
Sections~\ref{sec:dysthe}--\ref{sec:reading} beside the paper.
Appendix~\ref{sec:details} fills in two compressed derivative calculations.

You are ready for a first reading when you can explain
\begin{enumerate}
\item why the carrier oscillates rapidly but its envelope changes slowly;
\item why a deep-water Fourier mode decays exponentially below the surface;
\item why group velocity is half the phase velocity;
\item why a spatially varying intensity generates a mean current;
\item why a negative radicand in the stability formula means growth.
\end{enumerate}
For a derivation-level reading, also work through the determinant in
Section~\ref{sec:stability} and the order reduction in
Appendix~\ref{sec:details}. The exercises have answers.

\subsection{What is derived here and what is taken from the paper}
The linear water-wave dispersion relation, NLS perturbation calculation,
Fourier solution for the mean current, Dysthe stability determinant, and
small-steepness growth formulas are worked out here. The numerical
coefficients in the full fourth-order equation are transcribed from the
scanned paper. The guide explains the harmonic-expansion procedure that
produces them and the subsequent order reduction; it does not repeat
every nonlinear product in the intermediate quantity $Q$.

Agreement with Longuet-Higgins's full-equation calculations or laboratory
observations refers to \emph{comparisons reported in the 1979 paper}.
These are not new computations or independent validations. Our plots of
approximate formulas do not reproduce those reference data. Perturbative
derivations below are formal asymptotic calculations unless a proposition
explicitly states an exact result for a specified reduced model.

\subsection{Keep the different waves and frequencies separate}
The \emph{carrier} has dimensional wavenumber $k_0>0$ and linear angular
frequency $\omega_0=\sqrt{gk_0}$. The \emph{modulation} is a slower
variation of amplitude and phase. In the paper's dimensionless stability
calculation its wavenumbers are $\lambda,\mu$ and its frequency is $\Omega$.
Thus $\Omega$ is not $\omega_0$, and $\lambda$ is not the carrier wavelength.

\begin{longtable}{@{}p{.20\linewidth}p{.74\linewidth}@{}}
\toprule
Symbol & Meaning and convention\\
\midrule
\endhead
$x,y,z,t$ & Along-wave, transverse, vertical coordinates and time;
water lies below $z=\zeta$.\\
$g,\rho$ & Gravitational acceleration and constant water density.\\
$\zeta,\phi$ & Real surface elevation and real velocity potential;
fluid velocity is $\nabla\phi$.\\
$a,\varepsilon$ & Leading crest amplitude and steepness $\varepsilon=k_0a$;
wave height is approximately $2a$.\\
$\theta$ & Carrier phase $k_0x-\omega_0t$.\\
$A,B$ & Complex first-harmonic coefficients of potential and elevation.
They are different variables.\\
$\Ph,\Zh$ & Mean potential and mean elevation; the bar denotes a
carrier-period average.\\
$A^*,|A|^2$ & Complex conjugate and $AA^*$.\\
$A_0$ & Real dimensionless constant envelope amplitude; $A_0=k_0a$.\\
$a',\theta'$ & Small real amplitude and phase perturbations in (3.1);
primes here are not derivatives.\\
$\lambda,\mu,K$ & Dimensionless modulation wavenumbers in $x,y$ and
$K=\sqrt{\lambda^2+\mu^2}$.\\
$\Omega,\gamma$ & Modulation frequency and growth rate;
$\gamma=\operatorname{Im}\Omega$ for the growing mode.\\
$P_1,P_2;\ \Ps,\Pt$ & Differential operators in (3.2)--(3.3),
and algebraic quantities used in the dispersion relation.\\
\bottomrule
\end{longtable}

Coordinates are dimensional until Section~\ref{sec:scaling}, and thereafter
use Dysthe's dimensionless convention unless stated otherwise. Paper
equation numbers appear as \papeq{2.19}; other equation numbers refer to
this guide. Printed pages 105--114 correspond to PDF pages 2--11 because
the supplied file begins with a JSTOR cover page.
\clearpage
